%%%%%%%%%%%%%%%%%%%%%%%%%%%%%%%%%%%%%%%%%%%%%%%%%%%%%%%%%%%%%%%%%%%%%%%%%%%%%%%
%                                                                             %
% Anwendungsbeispiel für das Beamer-Template TU Chemnitz                      %
% (c) Mario Haustein (mario.haustein@hrz.tu-chemnitz.de), 2013-2019           %
%                                                                             %
%%%%%%%%%%%%%%%%%%%%%%%%%%%%%%%%%%%%%%%%%%%%%%%%%%%%%%%%%%%%%%%%%%%%%%%%%%%%%%%

\usepackage[T1]{fontenc}
\usepackage[utf8]{inputenc}
\usepackage[ngerman]{babel}


% TUC-Templates laden.
\usetheme{tuc2019}
\mode<article>{\usepackage{beamerarticletuc2019}}


%
% Weitere Anpassungen nach Bedarf.
%

%% Navigationsleiste deaktivieren.
%\setbeamertemplate{navigation symbols}{}

%% Mathematische Sätze nummerieren.
%\setbeamertemplate{theorems}[numbered]

%% Titelzeile fett und in Standardfarbe formatieren.
%\setbeamercolor{frametitle}{parent=normal text}
%\setbeamerfont{frametitle}{series=\bfseries}

% Metadaten
\newcommand{\Title}{Benutzerfreundlichkeit fünf verschiedener Zwei-Faktor-Authentifizierungsmethoden}
\newcommand{\ShortTitle}{Zwei-Faktor-Authentifizierungsmethoden}
\title[\ShortTitle]{\Title}
\subtitle{Hauptseminar Medieninformatik}
\author{Joshua Jeschek}
\date{11.06.2022}
\institute[TUC]{TU Chemnitz}
\titlegraphic{\includegraphics[height=0.19\textheight]{tuc2019/logo/tuc_green}}
\tucurl{https://hauptseminar.jeschek.eu}


\begin{document}
\tucthreeheadlines{}
\frame{\titlepage}

\tuctwoheadlines{}
\section*{\ShortTitle}
\subsection{Gliederung}
\frame{\frametitle{Gliederung}\tableofcontents[hideallsubsections]}

%
% Inhalt
%

\tucthreeheadlines{}
\section{}
\frame<beamer| trans>{\begin{center}\scalebox{2}{Vielen Dank für eure Aufmerksamkeit.}\end{center}}

% \begin{onlyenv}<beamer| trans>
% \appendix
% \newcounter{finalframe}
% \setcounter{finalframe}{\value{framenumber}}

% %
% % "Spare-Slides" hier einsetzen.
% %

% \setcounter{framenumber}{\value{finalframe}}
% \end{onlyenv}

\tuctwoheadlines{}
\section*{\ShortTitle}
\subsection{Literatur}
\begin{frame}
\frametitle{Literatur}

\begin{thebibliography}{xxx}
\bibitem{reese2019}
    Reese, K., Smith, T., Dutson, J., Armknecht, J., Cameron, J. \& Seamons, K., 2019 \---
    A Usability Study of Five Two-Factor Authentication Methods. SOUPS @ USENIX Security Symposium.
\newblock{\url{https://www.usenix.org/conference/soups2019/presentation/reese}}
% \newblock{}
\end{thebibliography}
\end{frame}


\end{document}
